% Cover letter — Hp3 manuscript, Cell Reports (Cell Press)
% Build:  latexmk -pdfxe cover_letter-cr.tex   (xelatex)
% Variant of cover_letter.tex (Microbial Biotechnology); differs in addressee,
% journal-fit framing (broad significance for Cell Reports), and submission line.
% Suggested-reviewer panel reused verbatim (all are phage-genomics / Aeromonas /
% fish-in-vivo experts equally appropriate for Cell Reports).
\documentclass[11pt,a4paper]{article}
\usepackage{iftex}
\ifPDFTeX\usepackage[T1]{fontenc}\usepackage[utf8]{inputenc}
\else\usepackage{fontspec}\setmainfont{Times New Roman}\fi
\usepackage[a4paper,margin=1in]{geometry}
\usepackage{parskip}
\usepackage{hyperref}\hypersetup{hidelinks}
\pagestyle{empty}

\begin{document}

\noindent Jinlong Ru, on behalf of all authors\\
Technical University of Munich \& Helmholtz Centre Munich\\
\href{mailto:jinlong.ru@tum.de}{jinlong.ru@tum.de}

\bigskip
\noindent The Editors\\
\emph{Cell Reports}

\bigskip
\noindent Dear Editors,

We are pleased to submit our manuscript, \textbf{``A novel lytic
\emph{Aeromonas} bacteriophage rescues largemouth bass from lethal
\emph{Aeromonas hydrophila} infection,''} for consideration as an Article in
\emph{Cell Reports}.

\emph{Aeromonas hydrophila}, a leading pathogen of global freshwater
aquaculture, is increasingly multidrug-resistant, eroding antibiotic control
and intensifying the search for sustainable alternatives such as phage therapy.
Our study connects three findings that we believe are of broad interest across
microbiology, viral evolution, and translational phage biology:

\begin{enumerate}
\item \textbf{A deeply divergent new viral genus.} Integrating nucleotide-level
(VIRIDIC, 83.65\% intergenomic similarity), whole-proteome (ViPTree),
amino-acid genome-BLAST distance phylogeny (VICTOR), and protein-sharing
network (vConTACT3) analyses, Hp3 and its closest relative HJ05 form a new,
deeply divergent genus within the class \emph{Caudoviricetes}---with no more
than $S_G = 0.073$ whole-proteome similarity to any classified virus and clear
separation from all established families. This places Hp3 in an
under-explored region of viral sequence space and indicates that the phages of
a globally important aquaculture pathogen remain substantially undersampled.

\item \textbf{Demonstrated \emph{in vivo} efficacy and host safety.} A single
administration (MOI 1.0) increased the seven-day survival of largemouth bass
from 15\% to 85\% under lethal challenge, while fish receiving phage alone
showed 100\% survival with no histopathology.

\item \textbf{A genomically safe, well-defined biocontrol candidate.} Hp3
carries no antibiotic-resistance genes and no high-confidence virulence
determinants, and exhibits rapid lytic kinetics with broad temperature and pH
stability.

\end{enumerate}

Together, these results link the discovery of a genomically distinct virus to a
measurable therapeutic outcome and frame Hp3 as a rigorously characterized,
genomically defined candidate for phage-based biocontrol---and as a tractable
starting point for the mechanistic and formulation studies that phage therapy
against \emph{A.~hydrophila} now requires. We believe the combination of a
deeply divergent new lineage, a defined genomic-safety profile, and an
\emph{in vivo} survival benefit will interest the broad readership of
\emph{Cell Reports}.

The manuscript is original, has not been published previously, and is not under
consideration elsewhere. As corresponding author, I consent to its publication
and confirm that the content and authorship of this manuscript have been
approved by all authors; the authors declare no competing interests. All raw
sequencing data and assembled genomes are openly available in NCBI BioProject
\textbf{PRJNA1348731} (Hp3 genome under GenBank \textbf{PX754746.1}); figure
source data are deposited in Figshare
(\href{https://doi.org/10.6084/m9.figshare.30500525}{10.6084/m9.figshare.30500525})
and the analysis code is archived in Zenodo
(\href{https://doi.org/10.5281/zenodo.17454899}{10.5281/zenodo.17454899}).
All animal experiments followed the ARRIVE guidelines and were approved by the
Institutional Animal Care and Use Committee (approval no. NWAFU-314020038).

% --------------------------------------------------------------------------
% Suggested reviewers — five experts covering the three pillars of the paper
% (phage genomics/taxonomy; Aeromonas/aquaculture phage therapy; fish in-vivo
% disease models). None are collaborators or co-authors of the authors.
% Substitute bench: Dann Turner (UWE Bristol, Dann2.Turner@uwe.ac.uk) for a
% pillar-1 decline.
% --------------------------------------------------------------------------
\medskip
\noindent\textbf{Suggested reviewers} (none has a conflict of interest; none is a
recent collaborator or co-author of the authors):
\begin{itemize}
\item \textbf{Evelien M. Adriaenssens}, Quadram Institute Bioscience, Norwich, UK
(\href{mailto:evelien.adriaenssens@quadram.ac.uk}{evelien.adriaenssens@quadram.ac.uk})
--- bacteriophage genomics and prokaryotic virus taxonomy; Chair of the ICTV
Bacterial Viruses Subcommittee (\emph{Caudoviricetes} classification framework).
\item \textbf{Andrew D. Millard}, University of Leicester, UK
(\href{mailto:adm39@leicester.ac.uk}{adm39@leicester.ac.uk}) --- bacteriophage
comparative genomics, reference-genome curation, and therapeutic phages.
\item \textbf{Mao Lin}, Fisheries College, Jimei University, Xiamen, China
(\href{mailto:linmao@jmu.edu.cn}{linmao@jmu.edu.cn}) --- \emph{Aeromonas} and
aquaculture bacteriophage genomics and phage therapy (host-range and genome
characterization).
\item \textbf{Jianguo Su}, College of Fisheries, Huazhong Agricultural University,
Wuhan, China (\href{mailto:sujianguo@mail.hzau.edu.cn}{sujianguo@mail.hzau.edu.cn})
--- fish antibacterial immunity and \emph{A.~hydrophila} \emph{in vivo} challenge
models.
\item \textbf{Lingbing Zeng}, Yangtze River Fisheries Research Institute, Chinese
Academy of Fishery Sciences, Wuhan, China
(\href{mailto:zenglingbing@gmail.com}{zenglingbing@gmail.com}) --- fish pathology
and largemouth bass (\emph{Micropterus salmoides}) disease models.
\end{itemize}

\medskip
\noindent We thank you for considering our work and look forward to your
response.

\bigskip
\noindent Sincerely,\\[1ex]
Jinlong Ru, Li Deng, and Gaoxue Wang\\
\emph{(corresponding authors)}\\
on behalf of all authors

\end{document}
